\scanpage{10}
\mspage{3}
\uline{Proof} $(i)\Rightarrow(ii)$ Assume (i), let $f\colon X\to Y$, $f\in F$, prove \deleted{$\uline{\mathrm{Hom}}(T,f)$}
\[
(\uline{\mathrm{Hom}}(T,f)\colon\uline{\mathrm{Hom}}(T,X)\to\uline{\mathrm{Hom}}(T,Y))\in F
\]
i.e. RLP with respect to any $i\colon A\to B$ \inserted{in $TC$} --- this amounts to RLP of $f$ with respect to $i\times\mathrm{id}_T\colon A\times T\to B\times T$, OK as this is \inserted{still} in $TC$.
Same for $f\in\unsure{TC|TF}$, using stability of $C$ under $\times T$.

$(ii)\Longrightarrow(i)$ Assume (ii), let $i\colon A\to B$, \deleted{prove} \inserted{in $C$}, prove $i\times T\colon A\times T\to B\times T$ is still in $C$, i.e. has \deleted{right} LLP with respect to any $f\in TF$, but this amounts to $i$ having the same with respect to $\uline{\mathrm{Hom}}(T,f)$, which is still in $TF$ by assumption OK.

\deleted{$\unsure{(i)\mathbin{\&}(ii)}\Longrightarrow(iii)$}
\uline{NB} The equivalence $(i)\Leftrightarrow(ii)$ can be further generalized, as a relation between $\unsure{\Phi^*},\unsure{\Phi_*}$ associated by LLP and RLP, when we have two functors $\otimes,\uline{\mathrm{Hom}}$ tied by the Cartan relation
\[
\mathrm{Hom}(A\otimes T,X)\simeq\mathrm{Hom}(A,\uline{\mathrm{Hom}}(T,X)).
\]
Moreover, when we have stability of $\unsure{\Phi^*}$ under \deleted{$\unsure{?\otimes?}\ T\otimes?$} \inserted{$?\otimes T$} \inserted{(where \unsure{$\Phi^*$ is} either $C$ or $TC$)}
\deleted{$\unsure{(i)\mathbin{\&}(ii)}\Longrightarrow\unsure{(iii)}$}
or equivalently of $\unsure{\Phi_*}$ under $\uline{\mathrm{Hom}}(T,?)$, then for $i\colon A\to B$ in $\Phi^*$, $f\colon X\to Y$ in $\unsure{\Phi_*}$, we'll have that
\[
\inserted{\uline{$\alpha(i,f)$}}\deleted{\uline{$\mathrm{Hom}(i,f)$}}\deleted{\unsure{$\alpha$}}\colon
\uline{\mathrm{Hom}}(B,X)\to\uline{\mathrm{Hom}}(i,f)
\quad\uline{\text{is in }TF}\inserted{(or in $\Phi_*?$)}.
\]
Indeed, this means it has RLP with respect to any $i'\colon A'\to B'$ in $C$ \inserted{[$\Phi^*$]}, i.e.
\[
\mathrm{Hom}(B',\uline{\mathrm{Hom}}(B,X))\to
\mathrm{Hom}(\unsure{i''|i'},\deleted{\unsure{\alpha(i,f)}}\inserted{\uline{\mathrm{Hom}}(\deleted{\unsure{\alpha}}\unsure{i},f)})
\]
is surjective. But \deleted{the first \unsure{hand}} \deleted{\unsure{this can}} \inserted{this can be viewed as}
\[
\deleted{\mathrm{Hom}(B'\otimes B,X)}
\]
\[
\mathrm{Hom}(B'\otimes B,X)\longrightarrow\mathrm{Hom}(\deleted{\unsure{\alpha(i',i)}}\deleted{\unsure{i'\otimes i}}\inserted{i'\boxtimes i},f)
\]
where
\[
\deleted{i'\otimes i\colon A'\otimes A\to B'\otimes B}
\]
\[
i'\boxtimes i\quad\deleted{\unsure{\alpha(i,i')}}\deleted{\unsure{i'\boxtimes i}}\colon\quad
A'\otimes B\amalg_{A'\otimes A}B'\otimes A\longrightarrow B'\otimes B.
\]
\margin{\unsure{Much longer}\ \unsure{more carefully}}
\deleted{and} Thus we have to check that if $i,i'\in\Phi^*$, then $i\boxtimes i'$ is in $\Phi^*$.
\margin{\unsure{and conversely}}
[OK if $\Phi^*$ are monomorphisms, say, and $\otimes={}$cartesian product]

Coming back to case of $(C,TF),(TC,F)$, this shows that if $i\in C$, $f\in TF$, then \deleted{$\unsure{\mathrm{Hom}}$} $\alpha(i,f)$ is in $TF$ --- and if $i$ is in $TC$ and $f$ in $F$ then \deleted{$\unsure{\mathrm{Hom}}$} $\alpha(i,f)$ is in $F$,
\mspage{9}

